(1) 一维Maxwell分布 \(f(v_x)=\sqrt{m/(2\pi k_BT)}e^{-mv_x^2/(2k_BT)}\)。变量替换 \(\omega=\omega_0+kv_x\) 给出 \(g(\omega)=[\sqrt{2\pi}\sigma_\omega]^{-1}\exp[-(\omega-\omega_0)^2/(2\sigma_\omega^2)]\)，其中 \(\sigma_\omega=|k|\sqrt{k_BT/m}\)。
(2) 对称性给出 \(\langle\omega\rangle=\omega_0\)，方差 \(\langle(\omega-\omega_0)^2\rangle=k^2k_BT/m\)。解 \(g(\omega_0\pm\Delta/2)=g(\omega_0)/2\)，得到半高全宽 \(\Delta\omega_{\rm FWHM}=2\sqrt{2\ln2}\,\sigma_\omega\)。
(3) 本征Lorentz谱 \(L_\Gamma(\omega)=\Gamma/[\pi(\omega^2+\Gamma^2)]\)，观测谱为Voigt卷积 \(V(\omega)=\int_{-\infty}^\infty g(\omega')L_\Gamma(\omega-\omega')\,d\omega'\)。由于两个核均归一，\(\int V(\omega)d\omega=1\)，谱线展宽不改变总归一强度。
(4) \(T\to0\) 时 \(g\) 趋向 \(\delta(\omega-\omega_0)\)，\(V\) 退化为Lorentz线；\(\Gamma\to0\) 时 \(V\to g\)。近似要求 \(v_{\rm th}/c\ll1\) 且 \(|\delta\omega|\ll\omega_0\)；高速气体必须使用相对论多普勒关系。
